\subsection{\textit{Sprint} 3}

Para a penúltima \textit{sprint}, o time havia previsto a integração de todas as telas
do \textit{front-end} com o \textit{back-end}, além de entregar os débitos técnicos que haviam ficado
da \textit{sprint} passada. Quanto ao meu \textit{squad}, ficamos com a integração da página de
administração que já havia sido iniciada, além de assumir a criação dos débitos
técnicos dos modais de edição e deleção de usuários.

Quanto aos débitos técnicos, o time inteiro foi capaz de entregar.
Especificamente da minha parte, os modais de edição e deleção de usuário foram
entregues com certa facilidade. Já as integrações foram uma dor de cabeça para
fazer, sendo perceptível que nenhuma página chegou a ser completamente
integrada. Pelo menos, é possível dizer que todas as integrações planejadas estão
bem alinhadas e não devem ser muito complicadas de finalizar.

Apesar de acreditar que a inércia de código seria um dos maiores problemas,
ela não foi fator, pelo menos não para o meu \textit{squad}. No entanto, foi possível ver
membros do time desanimados e partes do projeto andando mais devagar do que
deveriam. Tanto minha equipe quanto as outras do \textit{front-end} encontraram certas
complicações na integração, em parte por inexperiência, mas também em parte por rotas
não implementadas do \textit{back-end}.

Como integrante do time, percebi que dependo dos outros integrantes do time
e do trabalho deles. Sendo assim, mediante momentos complicados de desmotivação
do time, aprendi que é necessário dar um passo à frente e se preocupar com o
desenvolvimento dos outros companheiros de time, zelando pelo ânimo dos
companheiros e tentando incentivá-los sempre que possível para garantir que o time
esteja trabalhando de maneira uniforme.

Além disso, a persistência se mostrou uma grande aliada quanto à
resolução de muitos \textit{bugs}. Da parte técnica, aprendi fontes valiosas para obter
informações do projeto, como por exemplo, a documentação do Swagger
\cite{swagger-docs}. Também
tive muito contato com a aba de desenvolvedor do navegador, já que foi necessário
depurar a parte das requisições, além de tratar diversos códigos de erro.

Com a \textit{sprint} final se aproximando, as \textit{tasks} que sobraram são as pendências,
além de dar uma refinada no código visualmente e estruturalmente.
Por enquanto, recebi a \textit{task} de continuar com a integração da página de
administrador, além de adicionar um botão para baixar o relatório e fazer este botão funcionar. Com
a finalização dessas \textit{tasks}, falarei com outros membros da equipe e trabalharei em alguma das \textit{tasks} pendentes.
